% Options for packages loaded elsewhere
\PassOptionsToPackage{unicode}{hyperref}
\PassOptionsToPackage{hyphens}{url}
%
\documentclass[
  11pt,
]{article}
\usepackage{lmodern}
\usepackage{amssymb,amsmath}
\usepackage{ifxetex,ifluatex}
\ifnum 0\ifxetex 1\fi\ifluatex 1\fi=0 % if pdftex
  \usepackage[T1]{fontenc}
  \usepackage[utf8]{inputenc}
  \usepackage{textcomp} % provide euro and other symbols
\else % if luatex or xetex
  \usepackage{unicode-math}
  \defaultfontfeatures{Scale=MatchLowercase}
  \defaultfontfeatures[\rmfamily]{Ligatures=TeX,Scale=1}
\fi
% Use upquote if available, for straight quotes in verbatim environments
\IfFileExists{upquote.sty}{\usepackage{upquote}}{}
\IfFileExists{microtype.sty}{% use microtype if available
  \usepackage[]{microtype}
  \UseMicrotypeSet[protrusion]{basicmath} % disable protrusion for tt fonts
}{}
\makeatletter
\@ifundefined{KOMAClassName}{% if non-KOMA class
  \IfFileExists{parskip.sty}{%
    \usepackage{parskip}
  }{% else
    \setlength{\parindent}{0pt}
    \setlength{\parskip}{6pt plus 2pt minus 1pt}}
}{% if KOMA class
  \KOMAoptions{parskip=half}}
\makeatother
\usepackage{xcolor}
\IfFileExists{xurl.sty}{\usepackage{xurl}}{} % add URL line breaks if available
\IfFileExists{bookmark.sty}{\usepackage{bookmark}}{\usepackage{hyperref}}
\hypersetup{
  hidelinks,
  pdfcreator={LaTeX via pandoc}}
\urlstyle{same} % disable monospaced font for URLs
\usepackage[margin=1in]{geometry}
\setlength{\emergencystretch}{3em} % prevent overfull lines
\providecommand{\tightlist}{%
  \setlength{\itemsep}{0pt}\setlength{\parskip}{0pt}}
\setcounter{secnumdepth}{-\maxdimen} % remove section numbering
\usepackage{esint}
\setmathfont{Asana Math}
\setlength{\emergencystretch}{3em}

\author{}
\date{}

\begin{document}

\hypertarget{thin-fiber-monodromy-global-reduction-and-einstein-realization}{%
\section{Thin-fiber monodromy, global reduction, and Einstein
realization}\label{thin-fiber-monodromy-global-reduction-and-einstein-realization}}

24 September 2026. This note removes the reference-fiber thickness
hypothesis from the preceding cycle theorem, at the cost of an explicit
alternative: finite monodromy or a common invariant essential
multicurve. For irreducible monodromy the alternative gives exact
negative Kähler--Einstein existence on the same smooth total space. For
reducible monodromy it gives an actual decomposition of the surface
bundle, not an assertion that its pieces already have compatible
Einstein metrics.

The proof uses classical hyperbolic collars, quasiconformal length
distortion, and compactness of thick moduli. It does not invoke an
unverified Teichmüller Margulis theorem or claim novelty for these
classical facts. The project-level extension is their composition with
the actual cycle budget and pointwise Ricci input. The previous
\href{20260924_cycle_holonomy_realization.md}{cycle theorem} and
\href{20260924_ricci_monodromy_realization.md}{Ricci theorem} are
inputs. Their finite-monodromy realization is retained, not
rediscovered.

\hypertarget{statements-with-original-geometric-inputs}{%
\subsection{1. Statements with original geometric
inputs}\label{statements-with-original-geometric-inputs}}

Let \[
 F_f\longrightarrow M\xrightarrow{\pi}B_b,\qquad f,b\ge2,
\tag{1}
\] be a smooth oriented bundle of closed oriented surfaces. Its
monodromy is \(\rho:\pi_1(B,b_0)\to\operatorname{Mod}(F_f)\). A
multicurve means a finite nonempty set of distinct isotopy classes of
disjoint essential simple closed curves. Call the image of \(\rho\)
\textbf{irreducible} when it preserves no such multicurve setwise. This
is a property of the entire group; it is not a claim that each element
is pseudo-Anosov or that the image is infinite.

For any smooth metric \(g\) use its orthogonal Ehresmann connection. For
a loop \(\gamma\) based at \(b_0\), let \(\Phi_t\) denote actual
horizontal transport and set \[
 h_t=\Phi_t^*(g|_{F_{\gamma(t)}}),\qquad
 A_t=h_t^{-1}\dot h_t,\qquad
 A_t^0=A_t-\tfrac12\operatorname{tr}(A_t)I.
\] For a finite based generating system \(\Gamma\) define \[
 \mathcal J_g(\gamma)
   =\frac12\int_0^1\sup_{F_{b_0}}\|A_t^0\|_{\rm op}\,dt,
 \qquad
 \mathcal J_g(\Gamma)=\max_{\gamma\in\Gamma}\mathcal J_g(\gamma).
\tag{2}
\] All loops in \(\Gamma\) have the same base point. No representative
chosen merely to make the mapping-class displacement small replaces the
transport in (2).

\textbf{Theorem A (global thin-fiber alternative).} There is a constant
\(\mu_f>0\), depending only on fiber genus, such that \[
                         \mathcal J_g(\Gamma)<\mu_f             \tag{3}
\] has the following consequence. Either the monodromy image is finite,
or it preserves a multicurve \(C\) with at most \(3f-3\) components. In
the latter case \(C\) can be chosen so that every component has length
less than \(1\) in the curvature \(-1\) uniformization of the reference
fiber.

If the monodromy is irreducible, it therefore has finite image and the
same smooth \(M\) carries a metric \(\widehat g\) with \[
 \operatorname{Ric}_{\widehat g}=-\widehat g,\qquad
 (\widetilde M,\widetilde{\widehat g})
     \cong\mathbb H^2_{-1}\times\mathbb H^2_{-1}.              \tag{4}
\] The output is Kähler in both orientations. There is no reference
systole assumption, no prescribed Einstein background, and no claim that
the input metric is already Einstein or close to the output. The
constant is independent of base genus and generator count.

The alternatives need not be disjoint: a finite group can preserve a
multicurve. The assertion relevant to an infinite image is that one
\textbf{common} multicurve is preserved by every element, not merely
that each generator has its own reduction system.

\textbf{Theorem B (pointwise Ricci realization without a thick fiber).}
Suppose \(\pi:(M,g)\to(B,k)\) is a Riemannian submersion with minimal
fibers. Intrinsic fiber curvature is unrestricted and the horizontal
distribution need not be integrable. Define \[
 \delta=\|\operatorname{Ric}_g^{\sharp}+I\|_{\infty,\rm op},
 \quad \eta=\|(K_k+1)_+\|_\infty,
 \quad D_B=\operatorname{diam}(B,k).
\tag{5}
\] If \[
                         2D_B\sqrt{\delta+\eta}<\mu_f,          \tag{6}
\] then the finite-or-reducible conclusion of Theorem A holds. In
particular, if \(\rho\) is irreducible, (4) holds on the original smooth
total space. For a hyperbolic base \(K_k=-1\), this is the sufficient
condition \[
                   \delta<\frac{\mu_f^2}{4D_B^2}.             \tag{7}
\] Consequently every such submersion with \textbf{infinite irreducible}
monodromy satisfies the input obstruction \[
                   (\delta+\eta)D_B^2\ge\frac{\mu_f^2}{4}.   \tag{8}
\] This concerns the specified submersion metrics, not arbitrary metrics
on \(M\). In the irreducible small-input case it nevertheless proves the
exact Einstein existence requested by Yau26 for that input class. No
total-space diameter, volume, injectivity-radius or auxiliary monopole
assumption is used.

\hypertarget{a-simultaneous-short-curve-lemma}{%
\subsection{2. A simultaneous short-curve
lemma}\label{a-simultaneous-short-curve-lemma}}

On Teichmüller space use \(d_T(x,y)=\tfrac12\inf\log K\). For a curve
class \(c\), write \(\ell_x(c)\) for its hyperbolic geodesic length.

Two classical facts enter the proof. First, the collar theorem implies
that all simple geodesics of length at most \(1\) are pairwise disjoint.
Indeed the collar inequality for intersecting simple geodesics gives
\(\sinh(\ell(c)/2)\sinh(\ell(d)/2)>1\), whereas \(\sinh^2(1/2)<1\). Thus
there are at most \[
                              N_f=3f-3                         \tag{9}
\] such classes. Primary sources are Keen,
\href{https://doi.org/10.1515/9781400881642-021}{\emph{Collars on
Riemann surfaces}}, and Buser,
\href{https://doi.org/10.1007/BF01168048}{\emph{The collar theorem and
examples}}.

Second, Wolpert's quasiconformal length inequality gives \[
 e^{-2d_T(x,y)}\ell_x(c)
       \le\ell_y(c)\le e^{2d_T(x,y)}\ell_x(c).                 \tag{10}
\] See \href{https://annals.math.princeton.edu/1979/109-2/p03}{\emph{The
length spectra as moduli for compact Riemann surfaces}}. The needed
normalization can also be checked directly. The cyclic cover
corresponding to \(c\), with core length \(\ell\), is a conformal
annulus of modulus \(\pi/\ell\). A \(K\)-quasiconformal map lifts to the
corresponding cyclic covers and changes annular modulus by a factor
between \(K^{-1}\) and \(K\). Take the infimum over marked maps to
obtain (10). This argument applies to every essential simple class, not
just a chosen pants decomposition.

Put \[
 a_f=\frac{\log2}{N_f+1}=\frac{\log2}{3f-2},
 \qquad \nu_f=\frac{a_f}{4}.                                 \tag{11}
\]

\textbf{Lemma 1 (a common reduction, with no generator-count bound).}
Suppose \(\operatorname{sys}(x)<1/2\), and let \(S\) be any set of
mapping classes satisfying \[
                          d_T(x,sx)<\nu_f\quad(s\in S).        \tag{12}
\] Then \(\langle S\rangle\) preserves a nonempty multicurve all of
whose \(x\)-lengths are less than \(1\).

\textbf{Proof.} In the interval \([-\log2,0]\), list the logarithms of
the lengths in \([1/2,1]\). There are at most \(N_f\) entries, including
repetitions if distinct curves have equal lengths. Insert the two
endpoints. Among the resulting at most \(N_f+1\) gaps there is one
\((u,v)\) with \[
                            v-u\ge a_f.                       \tag{13}
\] There is no curve length with logarithm in this open interval. Let
\(t=\exp((u+v)/2)\) and define \[
                            C=\{c:\ell_x(c)<t\}.              \tag{14}
\] It is nonempty because the systole is less than \(1/2\), and it is a
multicurve by the collar theorem. Every member of \(C\) has length at
most \(e^u\). The action convention may invert \(s\) in (10), but
\(d_T(x,s^{-1}x)=d_T(x,sx)\), so \[
 \ell_x(sc)\le e^{2d_T(x,sx)}\ell_x(c)
     < e^{a_f/2}e^u\le e^{(u+v)/2}=t.                        \tag{15}
\] Hence \(sC\subset C\). Since \(s\) is a bijection of curve classes
and \(C\) is finite, \(sC=C\). The same \(C\) works for all generators,
and is therefore preserved by their entire group. Repeated length values
and curves lying at the endpoints of the gap cause no problem: the
strict inequality in (12) supplies the strict inequality in (15).
\(\square\)

This is a length-gap argument, not an assertion that small moves of the
generators make every group element a small move. Long words may
accumulate arbitrarily large displacement.

\textbf{Lemma 2 (thick finite-isotropy neighbourhood).} Some
\(\varepsilon_{f,1/2}>0\) has the property that, if
\(\operatorname{sys}(x)\ge1/2\) and \(d_T(x,sx)<\varepsilon_{f,1/2}\)
for every \(s\in S\), then \(\langle S\rangle\) is finite.

\textbf{Proof.} Mumford compactness gives a compact set in Teichmüller
space meeting every orbit with this systole bound. If the assertion
failed, conjugate to obtain \(x_j\to x\) in that compact set and sets
\(S_j\) generating infinite groups, with all displacements less than
\(1/j\). Uniformly in \(s\in S_j\), \[
                  d_T(x,sx)\le2d_T(x,x_j)+1/j\longrightarrow0.
\tag{16}
\] Properness of the mapping-class action and properness of the
finite-dimensional Teichmüller metric imply that the set of mapping
classes moving \(x\) by at most \(1\) is finite. Its positive
displacements have a positive minimum. Thus for large \(j\) every member
of \(S_j\) fixes \(x\). The stabilizer is finite, giving a
contradiction. A finite ineffective kernel, such as the genus-two
hyperelliptic involution, does not affect this argument. This is the
lemma proved in the preceding cycle note, with \(\iota=1/2\). The
classical compactness source is Mumford,
\href{https://doi.org/10.2307/2037802}{\emph{A remark on Mahler's
compactness theorem}}. \(\square\)

Fix once and for all \[
                \mu_f=\min\{\varepsilon_{f,1/2},\nu_f\}>0.   \tag{17}
\] Lemmas 1--2 give the following purely geometric alternative for every
\(x\in\mathcal T_f\): the subgroup generated by any collection of
\(\mu_f\)-small moves is finite or has a common short multicurve. The
second conclusion is not replaced by ``each generator is reducible.''

\hypertarget{from-the-actual-connection-to-the-group-alternative}{%
\subsection{3. From the actual connection to the group
alternative}\label{from-the-actual-connection-to-the-group-alternative}}

For completeness the differential-to-discrete step has no compactness
hypothesis. For two fixed nonzero tangent vectors on the reference
fiber, \[
 \left|\frac d{dt}\log\frac{h_t(v,v)}{h_t(w,w)}\right|
                      \le2\|A_t^0\|_{\rm op}.                \tag{18}
\] Integration bounds the logarithm of the quasiconformal dilatation of
the identity between \(h_0\) and \(h_1\) by
\(\int\sup\|A_t^0\|_{\rm op}\,dt\). The endpoint metric is the pullback
by the actual return map, which represents \(\rho(\gamma)\). Therefore
\[
                 d_T(x_0,\rho(\gamma)x_0)
                                   \le\mathcal J_g(\gamma).  \tag{19}
\] Noncommuting shape operators and a curved connection are allowed. The
mapping classes of return maps are invariant under based homotopy and
compose as monodromy, even though the actual return diffeomorphisms need
not form a representation.

Apply (17) to (19) for the generating system. This proves the
alternative in Theorem A. If \(\rho\) is irreducible, the short
multicurve is impossible, so the image \(G\) is finite. Kerckhoff's
\href{https://annals.math.princeton.edu/1983/117-2/p02}{Nielsen
realization theorem} realizes \(G\) by orientation-preserving isometries
of one hyperbolic fiber. On the connected kernel cover
\(\widehat B\to B\), use a lifted hyperbolic base metric. The diagonal
action on \(\widehat B\times F_f\) is free, regardless of fixed points
of the fiber action, since its base action is free. The product metric
descends with Ricci tensor equal to minus the metric.

This construction gives the original smooth bundle, not just a manifold
with an isomorphic group. For \(f\ge2\), Earle--Eells contractibility of
\(\operatorname{Diff}_0^+(F_f)\) gives
\(B\operatorname{Diff}^+(F_f)\simeq B\operatorname{Mod}(F_f)\). The two
bundles with the same monodromy are therefore smoothly
bundle-isomorphic. Pull back the product-quotient metric. The product
complex structures with either sign on the fiber descend, giving (4) in
both orientations. This concludes Theorem A. The primary contractibility
source is Earle--Eells,
\href{https://doi.org/10.1090/S0002-9904-1967-11746-4}{\emph{The
diffeomorphism group of a compact Riemann surface}}.

\hypertarget{ricci-input-and-the-scale-of-irreducible-monodromy}{%
\subsection{4. Ricci input and the scale of irreducible
monodromy}\label{ricci-input-and-the-scale-of-irreducible-monodromy}}

For a horizontal vector \(X\), write \[
 g(S_XU,V)=\langle(\nabla_UV)^{\mathcal H},X\rangle,
 \qquad A_XY=\tfrac12[X,Y]^{\mathcal V}.
\] Minimality gives \(\operatorname{tr}S_X=0\). The classical horizontal
Ricci identity, already used in retained Round 20, is \[
 \operatorname{Ric}_g(X,X)
      =K_k-|S_X|_{\rm HS}^2-2\sum_a|A_Xe_a|^2
                      \quad(|X|=1).                         \tag{20}
\] The mean-curvature derivative term vanishes because the fibers are
pointwise minimal. The source is O'Neill,
\href{https://doi.org/10.1307/mmj/1028999604}{\emph{The fundamental
equations of a submersion}}. Equation (5) consequently implies \[
                         |S_X|_{\rm op}
                               \le\sqrt{\delta+\eta}\,|X|.  \tag{21}
\] Horizontal transport differentiates the vertical metric by
\(h_t^{-1}\dot h_t=-2\Phi_t^*S_{\dot\gamma^{\mathcal H}}\). Thus (2)
gives \[
                    \mathcal J_g(\gamma)
                          \le\sqrt{\delta+\eta}\,
                                      \operatorname{Length}_k(\gamma).
\tag{22}
\]

For each \(a>0\), a finite based generating system of a closed
Riemannian manifold can be chosen with lengths below
\(2\operatorname{diam}+a\). One direct proof is to subdivide any loop
into sufficiently short arcs and join the subdivision points to the base
point by paths of length at most the diameter. Each resulting loop has
length below \(2\operatorname{diam}+a\). The set of their possible
homotopy classes is finite: in the proper universal cover it is
represented by deck transformations whose orbit points lie in a fixed
compact ball. These classes generate the fundamental group. No bound on
their number is needed here.

Under (6), choose \(a>0\) small enough that
\((2D_B+a)\sqrt{\delta+\eta}<\mu_f\) and apply Theorem A to (22). If
\(\delta+\eta=0\), every cycle budget is zero and the same argument
applies directly. This proves Theorem B and its contrapositive (8).

There is also an explicit thickness consequence, without the implicit
constant in Mumford compactness. If \(\rho\) is irreducible and
\emph{some} fiber has conformal hyperbolic systole less than \(1/2\),
base the short generators there. Lemma 1 and (22) give \[
 (\delta+\eta)D_B^2
                  \ge\frac{(\log2)^2}{64(3f-2)^2}.            \tag{23}
\] Hence, below this scale, irreducibility forces every fiber to be
conformally thick. This is not a lower bound for four-dimensional ball
volumes. It identifies the missing information in the earlier
thick-fiber theorem: here a global topological property supplies the
fiber thickness needed for finite-monodromy detection.

\hypertarget{the-reducible-alternative-is-a-genuine-global-decomposition}{%
\subsection{5. The reducible alternative is a genuine global
decomposition}\label{the-reducible-alternative-is-a-genuine-global-decomposition}}

\textbf{Proposition 3 (fiberwise realization of the invariant
multicurve).} If \(\rho\) preserves a multicurve \(C\) of size \(r\),
then \(M\) contains a smooth submanifold \(\mathcal C\) whose projection
to \(B\) is a bundle with fiber the corresponding union of \(r\)
circles. After a connected finite base cover of degree at most
\(2^r r!\), its circle components and orientations can be labeled
globally. Cutting the pulled-back total space along them gives genuine
surface bundles with boundary over that base cover.

\textbf{Proof.} Isotopy classes of curves in the fibers form a discrete
local system. The monodromy-invariant set \(C\) gives a global section
of the associated local system of multicurves. Uniformize the vertical
metric in each fiber to curvature \(-1\). These metrics vary smoothly
with the base: the conformal curvature equation has an invertible
linearization at its negative-curvature solution, and uniqueness makes
the local solutions agree.

In each fiber take the unique geodesic representatives of the specified
classes. In a local bundle chart the classes are fixed. Their
unparametrized geodesic representatives depend smoothly on the
hyperbolic metric: the periodic normal Jacobi operator along a closed
hyperbolic geodesic has no kernel. The representatives are disjoint, and
uniqueness makes the local families agree under changes of chart. Their
union is the required embedded smooth \(\mathcal C\). It is compact and
its projection is a submersion, so it is a bundle by local triviality
for proper submersions.

The action on oriented components maps \(\pi_1(B)\) to the signed
permutation group of order \(2^r r!\). Pass to the connected cover
associated with its kernel. An orientation of a curve, together with the
fiber orientation, orients its normal line. Thus every boundary side
obtained by cutting is labeled. Every component of a cut connected
surface meets some such boundary side; consequently the components are
also fixed. Fiberwise tubular neighborhoods and cutting now produce the
stated smooth bundles with boundary. No assertion that the circle
bundles are trivial is needed or made. \(\square\)

This decomposition keeps rather than discards residual monodromy. After
the finite cover, let \(S_v\) be the cut fiber components and let
\(\overline S_v\) be the punctured surfaces obtained by capping their
boundary circles with marked punctures. Restriction defines \[
 \rho_v:\pi_1(\overline B)\longrightarrow
                              \operatorname{PMod}(\overline S_v).
\tag{24}
\] For the subgroup fixing curves and their sides, the kernel of the
product restriction map is precisely the free abelian group \[
                         \mathcal T_C\cong\mathbb Z^r        \tag{25}
\] of Dehn twists about \(C\). A representative trivial after capping
can be isotoped to the identity off annular neighborhoods; its remaining
annular mapping classes are integer twists. Conversely all those twists
vanish after capping, and twists in distinct essential nonparallel
curves are independent. This is the standard multicurve restriction
sequence; see Hubbard--Koch, Proposition 2.5 of
\href{https://arxiv.org/abs/1301.0062}{\emph{An analytic construction of
the Deligne--Mumford compactification of the moduli space of curves}}.

Thus the remaining data are the representations (24) \textbf{and} the
annular twist data (25). If the joint image in (24) is finite, a further
finite base cover reduces the image to \(\mathcal T_C\). Theorems A--B
do not prove that additional premise. A common multicurve alone cannot
replace it: a mapping class supported on a once-punctured torus with
homological action
\(\left(\begin{smallmatrix}2&1\\1&1\end{smallmatrix}\right)\) is
pseudo-Anosov on that component and extends by the identity across its
boundary. It preserves the boundary curve but is not virtually a
multitwist. This example explains the retained restriction data; it is
not claimed to satisfy the smallness constant in (3).

In particular, neither this reduction nor a local bidisk description
solves the Einstein matching problem along \(\mathcal C\). No invariant
multicurve is silently promoted to a global marked product or to finite
monodromy.

\hypertarget{actual-tests-of-both-alternatives}{%
\subsection{6. Actual tests of both
alternatives}\label{actual-tests-of-both-alternatives}}

\hypertarget{a-nonempty-finite-irreducible-realization-class}{%
\subsubsection{6.1 A nonempty finite irreducible realization
class}\label{a-nonempty-finite-irreducible-realization-class}}

Finite image is compatible with irreducibility. On the smooth genus-two
curve \[
                           y^2=x^5-1
\tag{26}
\] consider \(\phi(x,y)=(\zeta_5x,y)\). Its action on holomorphic
one-forms \(dx/y,x\,dx/y\) has eigenvalues \(\zeta_5,\zeta_5^2\), and
its action on \(H^1(F;\mathbb Q)\) has the irreducible characteristic
polynomial \(\Phi_5(t)=t^4+t^3+t^2+t+1\).

If a multicurve were preserved, its at most three components would each
be fixed by this order-five action. A nonseparating component would give
an invariant rational line in \(H_1\), impossible by irreducibility of
\(\Phi_5\). A separating component divides the genus-two surface into
two one-holed tori. The order-five action cannot exchange these two
sides, and their homology would give a proper invariant rational
two-plane, also impossible. Thus \(\langle\phi\rangle\) is irreducible.

Map a surface group onto this cyclic group, take its kernel cover, and
form the diagonal product quotient. It satisfies (1) with nontrivial
finite irreducible monodromy. The quotient product metric has
\(\delta=\eta=\mathcal J=0\). Small, base-dependent vertical metric
perturbations with fixed fiber area form, supported over a base chart,
give genuinely nonproduct minimal submersion inputs with \(\delta\to0\).
For example, exponentiate a smooth trace-free self-adjoint vertical
endomorphism with a small base-dependent coefficient, leaving the
horizontal and base metrics fixed. The determinant and hence vertical
area form stay fixed; the full Ricci tensor varies smoothly. These
inputs need not be Einstein, but the conclusion still concerns the same
smooth bundle.

\hypertarget{the-thin-dehn-twist-family-realizes-the-other-branch}{%
\subsubsection{6.2 The thin Dehn-twist family realizes the other
branch}\label{the-thin-dehn-twist-family-realizes-the-other-branch}}

The retained fixed bundle \(E_D\), with monodromy
\(\rho(\gamma)=D^{u(\gamma)}\) for one nonseparating Dehn twist,
provides an actual infinite reducible example. Section 6 of the
preceding cycle note gives, on a fixed hyperbolic base and with minimal
intrinsically hyperbolic fibers, metrics \(g_\lambda\) such that \[
 \delta_-(g_\lambda)=O(\lambda),\qquad
 \mathcal J_{g_\lambda}(\Gamma)=O(\lambda),\qquad
              \operatorname{sys}(x_\lambda)\le\lambda\to0.
\tag{27}
\] The full Ricci matrix, not merely the horizontal block, was computed
there. These metrics also have actual opposite-orientation monopole
classes and bilateral near-saturation; they are not formal
almost-parallel tensor examples.

Here the common invariant curve is precisely available as the twisting
curve. The monodromy remains infinite cyclic however small the budget
becomes. With fiber genus \(f\) and base genus \(b\), the split
fundamental-group extension and the rank-one transvection on fiber
homology give \[
                              b_1(E_D)=2f+2b-1.               \tag{28}
\] It is odd, so \(E_D\) admits no Kähler structure. Consequently one
cannot replace the finite-or-reducible alternative by KE realization for
every small-budget bundle. This is not a claim that \(E_D\) has no
general Einstein metric. The retained example is used as a test, not
counted as a new construction.

\hypertarget{a-signature-zero-infinite-irreducible-test}{%
\subsubsection{6.3 A signature-zero infinite irreducible
test}\label{a-signature-zero-infinite-irreducible-test}}

Choose a pants decomposition \(P\) of \(F_f\) and, for each component
\(p\), a simple transverse curve \(q_p\) supported in the adjacent
one-holed torus or four-holed sphere. The collection \(P\cup\{q_p\}\)
fills: a curve disjoint from \(P\) is isotopic to a pants curve, and
that class meets its transverse curve. Let \(G\) be generated by the
twists in this collection.

The group \(G\) is infinite and irreducible. Indeed, if it preserved a
finite multicurve \(C\), a power of each generating twist would fix
every member \(c\) of \(C\). The annular intersection formula \[
                    i(T_a^n c,c)=|n|\,i(a,c)^2               \tag{29}
\] forces \(i(a,c)=0\) for every generating curve \(a\). The filling
property rules out such an essential \(c\). Infinitude already follows
from one nontrivial twist.

Realize this representation through a free quotient of a base surface
group: use one free generator for each twist, send one member of each
standard handle pair to that generator and the other to the identity.
Actual twist diffeomorphisms give a suspension bundle. The
representation extends over a handlebody, so the total four-manifold
bounds the associated oriented five-dimensional bundle and has signature
zero.

The inherited free-quotient construction in
\href{../mathematics/current/round20_construction.md}{Round 20, Section
2} already constructs \(\delta_-\to0\) metrics of the required minimal
hyperbolic-base submersion type on these fixed bundles. No such family
is new here. What (8) newly supplies in the project is the
thickness-independent bound \[
             \operatorname{diam}(B,k_j)
                         \ge\frac{\mu_f}{2\sqrt{\delta_-(g_j)}}.
\tag{30}
\] It remains valid even if one tries to degenerate the conformal fibers
simultaneously. The earlier signature barrier is zero on these objects;
the earlier cycle theorem required a thick reference fiber.
Irreducibility detects a global obstruction to small-cycle deformation
that both numerical signature and uncontrolled fiber pinching miss.
Equation (30) is an obstruction to a bounded-base-scale input, not to
Einstein existence on the underlying total manifold.

\hypertarget{what-has-changed-in-the-first-line-realization-mechanism}{%
\subsection{7. What has changed in the first-line realization
mechanism}\label{what-has-changed-in-the-first-line-realization-mechanism}}

The previous sufficient condition used a conformally thick fiber to
convert small cycle deformation into finite monodromy. The new argument
distinguishes two ways of escaping that step. In thin moduli, small
simultaneous displacement forces a common short multicurve; in thick
moduli it forces finite image. Thus global irreducibility supplies the
required thickness rather than assuming it. Under the submersion
hypotheses this is a genuine Ricci-input-to-Einstein-output theorem,
with a genus-only constant and no fiber systole in the statement.

The reducible case produces a smooth global cutting locus and retains
both component and annular monodromy. It is the specific alternative
exhibited by the actual \(E_D\) near-Einstein family. A further theorem
controlling those remaining transition data would be needed for an
Einstein realization there. The current conclusion neither presumes that
such a theorem holds nor treats the existence problem as settled by
local products.

The substantive inputs and their roles are separate: hyperbolic collars
and quasiconformal length distortion give Lemma 1; Mumford compactness
gives Lemma 2; O'Neill gives (20); Nielsen realization and
contractibility of \(\operatorname{Diff}_0\) give exact output once
finite monodromy has actually been detected. No new gauge solvability or
compactness theorem is asserted by this composition.

\end{document}
